\documentclass[11pt]{article}

\usepackage{amsmath,amssymb}
\usepackage{xurl}
\usepackage[colorlinks=true,linkcolor=blue,citecolor=blue,urlcolor=blue]{hyperref}

% Shared commands for notes in docs/paper-gaps/.

\DeclareUnicodeCharacter{2080}{\ifmmode{}_{0}\else\textsubscript{0}\fi}
\DeclareUnicodeCharacter{2081}{\ifmmode{}_{1}\else\textsubscript{1}\fi}
\DeclareUnicodeCharacter{2082}{\ifmmode{}_{2}\else\textsubscript{2}\fi}
\DeclareUnicodeCharacter{2099}{\ifmmode{}_{n}\else\textsubscript{n}\fi}

\newcommand{\Lean}{\textsc{Lean}}
\ExplSyntaxOn
\NewDocumentCommand{\leanid}{m}{
  \ifmmode
    \textsf{\detokenize{#1}}
  \else
    \group_begin:
    \urlstyle{sf}
    \tl_set:Nn \l_tmpa_tl {#1}
    \tl_replace_all:Nnn \l_tmpa_tl {\_} {_}
    \exp_args:NV \path \l_tmpa_tl
    \group_end:
  \fi
}
\ExplSyntaxOff
\providecommand{\tr}{\operatorname{tr}}
\newcommand{\tnleanIssueBase}{https://github.com/LionSR/TNLean/issues/}
\newcommand{\tnleanPrBase}{https://github.com/LionSR/TNLean/pull/}
\newcommand{\tnleanDocBase}{https://sirui-lu.com/TNLean/docs/}
\newcommand{\ghissue}[1]{\href{\tnleanIssueBase#1}{GitHub issue \##1}}
\newcommand{\ghpr}[1]{\href{\tnleanPrBase#1}{PR \##1}}
\newcommand{\doclink}[1]{\href{\tnleanDocBase#1.html}{\path{#1}}}

% Dirac notation and the identity operator, as in blueprint/src/macros/common.tex.
% \providecommand keeps note-local definitions authoritative.
\usepackage{dsfont}
\providecommand{\ket}[1]{|#1\rangle}
\providecommand{\bra}[1]{\langle #1|}
\providecommand{\braket}[2]{\langle #1 | #2 \rangle}
\providecommand{\ketbra}[2]{|#1\rangle\!\langle #2|}
\providecommand{\Id}{\mathds{1}}
\providecommand{\one}{\mathds{1}}
\providecommand{\C}{\mathbb{C}}
\providecommand{\MN}[1]{M_{#1}(\C)}
\providecommand{\MD}{M_D(\C)}

% Verdict marker: \gapnote{<kind>}{<status>}. Kind is the policy
% classification (clarification, local-correction, scope-restriction,
% unfaithful, false-source, open-gap); status is open, wip (elimination
% underway), resolved, or historical. Machine-read by the site index;
% prints a verdict line.
\newcommand{\gapnote}[2]{\par\noindent\textbf{Verdict:} #1 (#2).\par}


\title{Finite-Group Fixed-Point Realization in the One-Dimensional SPT Classification}
\author{}
\date{2026-09-26}

\begin{document}
\maketitle
\gapnote{scope-restriction}{open}

\section{Source statement}

Schuch, P\'erez-Garc\'ia, and Cirac, arXiv:1010.3732, Section~II.F.2,
classify symmetric injective matrix product states by the cohomology class of
their virtual projective representations. Their equality-of-phases argument
first replaces each state by its isometric form, then interpolates between
two maximally entangled bonds in a common direct-sum virtual space. The
interpolating parent Hamiltonians are commuting and remain gapped. The source
statement permits distinct endpoint physical representations, embedded in a
larger common on-site representation.

The assertion covers more than the existence of a tensor with a prescribed
virtual cocycle: it compares arbitrary symmetric injective endpoints by a
symmetric gapped path and proves that distinct classes cannot be joined by a
smooth MPS path.

\section{Formalized finite-group construction}

Let $G$ be a finite group, and let
$\omega:G\times G\to\mathbb C^\times$ be a genuine multiplicative
$2$-cocycle. The formalization takes the unit-modulus part of $\omega$, uses
its inverse as input to the cocycle-twisted left regular matrices, and thereby
obtains a unitary projective virtual representation with factor system in
$[\omega]$. Inverting the input is necessary because the chosen left regular
matrices have factor system equal to the inverse of their input phase.

On a bond space of dimension $D=|G|>0$, the local matrix-unit tensor has
physical dimension $D^2$ and letters $A^{(i,j)}=|i\rangle\langle j|$. Its
letters span the full matrix algebra. Conjugation by the virtual matrices
induces a linear unitary on-site representation on the physical matrix space;
the tensor is symmetric under this representation. The row convention in the
Lean definition is the transpose of the inverse adjoint action, so that the
physical twist of a letter is exactly its virtual conjugate.

The principal declaration is
\leanid{MPSTensor.exists_unitary_injective_fixedPoint_for_cocycle} in
\doclink{TNLean/MPS/Symmetry/MatrixUnitFixedPoint}. Its conclusion gives an
injective symmetric tensor, a unitary physical representation chosen from
$\omega$, and a virtual factor system cohomologous to $\omega$. For finite
$G$, the comparison of $H^2(G,\mathbb C^\times)$ with
$H^2(G,\mathrm U(1))$ is proved separately in
\leanid{TNLean.Algebra.circleH2EquivComplexUnitsH2}.

\section{Restriction and elimination plan}

The construction assumes a finite group and chooses both the physical Hilbert
space and its on-site action from the prescribed class. It does not address a
preassigned physical representation, arbitrary groups, or the deformation of
an arbitrary endpoint tensor to this fixed point. In particular, it does not
supply the continuous symmetric path, uniformly gapped parent Hamiltonians,
or path invariance needed for the classification theorem in
\path{blueprint/src/chapter/ch12_symmetry_spt.tex}.

The direct-sum virtual action, normalized interpolating bond, and unitary
on-site action are now formalized.  In independent-bond coordinates,
\leanid{MPSTensor.normalizedBondInterpolation_parent_gap_one} proves a
chain-length- and parameter-independent gap of at least one, while
\leanid{MPSTensor.mpv_weightedMatrixUnitInterpolation_eq_incomingBondProduct}
identifies the corresponding positive-length matrix product coefficients
after regrouping physical registers.  These are ingredients of the forward
construction, not a path theorem for arbitrary endpoints.

Every one-site injective tensor with exact unitary on-site symmetry now
admits a symmetric uniformly gapped parent path to an isometric physical
map. The declaration is
\leanid{MPSTensor.exists_symmetric_injective_parent_path_to_isometry}.
First, scalar normalization and a virtual gauge preserve all parent
interactions and put the tensor in canonical form. The virtual symmetry
can then be chosen unitary. The simultaneous equivariant polar
factorization gives a continuous injective tensor path.

The uniform gap is proved independently of the polar factorization.
For each injective tensor, an open-chain comparison provides a strict
Knabe window. Continuity of the finite-window ground projection and
Hamiltonian preserves a smaller bound locally; a finite subcover gives
one bound on a compact parameter set for all sufficiently long rings.
Continuity of the periodic ground-state projection handles the finitely
many shorter lengths. The resulting declaration
\leanid{MPSTensor.exists_uniform_parentHamiltonianES_gap_of_compact_isInjective_all_lengths}
has no additional gap hypothesis and covers every ring of at least two
sites. Canonical parent interactions are projections of norm at most one,
and tensor covariance implies symmetry of every periodic Hamiltonian.

One must also connect the independent-bond
Hamiltonian to the specified embedded endpoint parent Hamiltonians with
the required locality and smoothness, and establish that the resulting
path meets the chain-length-uniform condition of
\leanid{MPSTensor.SymmetricGappedInteractionPath}.  For the reverse implication,
\leanid{MPSTensor.unitary_projectivePath_factor_endpoints_cohomologous}
now proves cohomology invariance for a given continuous, fixed-positive-
dimension virtual projective path, for any group.  It remains to extract
such a virtual path from an arbitrary smooth symmetric MPS/parent-Hamiltonian
path, including the blocking and gauge choices in the source.  The source classification entry
remains unformalized until these steps are proved with its stated
hypotheses.

\end{document}
